
%%%===========================================

% ===
\section{Final Remarks}
% ===

This thesis aims to address a foundational limitation in mainstream economic theory: the persistent treatment of the environment as an exogenous constraint, external to the logic of production and disconnected from the behavior of economic agents. By proposing the endogenization of environmental resources into the structure of economic reasoning, the research seeks to integrate ecological boundaries into the decision-making processes of producers, firms, and institutions. It would explore the hypothesis that environmental preservation could become a rational, utility-maximizing strategy when ecological limits are clearly defined and economic incentives are structured accordingly.

The research would begin by investigating the theoretical origins and institutional consequences of treating nature as a non-scarce or marginal input. It would argue that this framing contributes to inefficiencies, externalities, and ecological collapse, especially when resource exploitation exceeds scientifically defined thresholds. Although instruments such as carbon markets, environmental certificates, and payments for ecosystem services already exist, they tend to operate as fragmented compensations or marginal rewards, rather than as central components of the economic system. The thesis would therefore aim to reinterpret these tools through a conceptual framework in which ecological rationality is aligned with economic incentives.

From this perspective, the second part of the thesis would seek to translate the proposed theoretical insights into a coherent structure capable of clarifying how different incentive configurations — whether economic, institutional, or territorial — might shape the behavior of agents facing environmental trade-offs. The core question would not be whether preservation is desirable in normative terms, but whether it could become the most advantageous course of action given the constraints and opportunities defined by the system. This would require situating the environment not as a passive boundary condition, but as an active variable influencing utility, productivity, and long-term resilience.

The final section of the research would aim to test whether the hypotheses outlined in the earlier chapters are supported by empirical patterns and institutional dynamics observed in a real-world context. The state of Mato Grosso — marked by high agricultural output and ecological diversity — would provide a critical setting in which to examine the feasibility of applying such a framework. Ultimately, the thesis would propose that an economy structured around commensal logic could enable the state to act not as a punitive enforcer, but as a catalyst of value — rewarding those who operate in synergy with environmental constraints, and fostering a productive system in which conservation is no longer a cost, but a dominant strategy.

This research introduces an integrated framework where environmental limitations are viewed not as externalities but as internal factors influencing productive decisions, potentially offering both theoretical and practical insights. On a theoretical level, it aims to challenge the prevailing paradigm by altering the principles of production theory, introducing ecological rationality into the agents' utility, and reframing conservation as a strategic option, not a limitation. Practically, it could guide the creation of localized policies and incentives that bolster both environmental resilience and economic success. By doing so, this work may create new possibilities for regional governance, sustainable certification schemes, and data-driven public planning, especially in areas like Mato Grosso where ecological and economic limits intersect. More generally, it proposes a paradigm in which environmental preservation is not in opposition to economic development but is seen as the most logical and beneficial path forward.


\newpage

% fim %